%%% thesis_ukr.tex %%%%%%%%%%%%%%%%%%%%%%%%%%%%%%%%%%%%%%%%%%%%%%%%
%%
%%%%%%%%%%%%%%%%%%%%%%%%%%%%%%%%%%%%%%%%%%%%%%%%%%%%%%%%%%%%%%%%%%%%%%%%
\documentclass[10pt,reqno]{amsart}
\usepackage{cmap}
\usepackage[T2A]{fontenc}
\usepackage[utf8]{inputenc}
\usepackage[ukrainian]{babel}
\usepackage{indentfirst,textcase}
\usepackage[a5paper,margin={15mm,15mm}]{geometry}

%%% Підключайте необхідні пакети тут %%%%%%%%%%%%%%%%%%%%%%%%%%%%%%%%%%%
\usepackage{graphicx}
%%% Помістіть свої локальні означення тут %%%%%%%%%%%%%%%%%%%%%%%%%%%%%%
\newcommand{\R}{\mathbb{R}}
\newcommand{\1}{1\mkern-4.5mu\mathrm{l}}
\newcommand{\set}[1]{\left\{#1\right\}}
\newcommand{\norm}[2][]{\left\lVert#2\right\rVert_{#1}}
\DeclareMathOperator{\Prob}{\mathsf{P}}

%%% Означення теоремоподібних оточень %%%%%%%%%%%%%%%%%%%%%%%%%%%%%%%%%%
\theoremstyle{plain}
\newtheorem{theorem}{Теорема}
\newtheorem*{theorem*}{Теорема}
\newtheorem{corollary}{Наслідок}
\newtheorem*{corollary*}{Наслідок}
\newtheorem{lemma}{Лема}
\newtheorem*{lemma*}{Лема}
\newtheorem{proposition}{Твердження}
\newtheorem*{proposition*}{Твердження}
\newtheorem{conjecture}{Гіпотеза}
\newtheorem*{conjecture*}{Гіпотеза}
\theoremstyle{definition}
\newtheorem{definition}{Означення}
\newtheorem*{definition*}{Означення}
\newtheorem{example}{Приклад}
\newtheorem*{example*}{Приклад}
\theoremstyle{remark}
\newtheorem{remark}{Зауваження}
\newtheorem*{remark*}{Зауваження}

%%%%%%%%%%%%%%%%%%%%%%%%%%%%%%%%%%%%%%%%%%%%%%%%%%%%%%%%%%%%%%%%%%%%%%%%

\begin{document}

\title{Нетривіальні узагальнення глибини Т'юкі для класів розподілів та аналіз рандомізованих апроксимацій}

%%% Інформація про автора
\author{В.Є.~Панаско}
\address{КПІ ім. Ігоря Сікорського, Київ, Україна}
\email{panasko.vitalii@lll.kpi.ua}

%% Не змінюйте ці два рядки!
\makeatletter\def\and{\unskip,\;\space\ignorespaces}\def\author@andify{\nxandlist{\and}{\and}{\and}}\makeatother
\maketitle\pagestyle{empty}\thispagestyle{empty}

\textbf{Актуальність.} Глибина Т'юкі є ключовою мірою центральності в багатовимірному аналізі. Точне аналітичне обчислення є складним для багатьох розподілів, що стимулює розвиток теоретичних узагальнень та ефективних чисельних методів, наприклад для однорідного розподілу на одиничній кулі та для стабільної сім'ї розподілів.

\medskip

\textbf{Мета.} Ми розширюємо визначення глибини Т'юкі для двох нетривіальних класів: (i) рівномірний (ізотропний) розподіл на кулі $B^p\subset\R^p$, та (ii) сімейство симетричних $\alpha$-стабільних розподілів. Крім того, аналізуємо рандомізовані алгоритми апроксимації для обчислення глибини у вищих вимірах.

\medskip

\textbf{Основні результати.} У випадку рівномірного розподілу на одиничній кулі ми довели, що для $\theta\in B^p$
\[
\mathrm{depth}(\theta)
=\frac12\bigl[1 - I_{\|\theta\|^2}\bigl(\tfrac12,\tfrac{p+1}2\bigr)\bigr],
\]
де $I_z(a,b)$ — регуляризована неповна бета-функція. Для симетричних $\alpha$-стабільних законів у $\R^2$ показали, що оптимальний поріг для напряму задовольняє
\[
c_{\max}(u,v)=
\begin{cases}
(u^\beta+v^\beta)^{1/\beta}, &\alpha>1,\\
\max\{u,v\}, &0<\alpha\le1,
\end{cases}
\]
з $\beta=\alpha/(\alpha-1)$. За допомогою чисельних експериментів було знайдено контури Тукі для біваріантних $t$- та експоненційних розподілів, показавши вплив параметрів на геометричний вигляд контурів.

\medskip

\textbf{Подальша робота.} На основі теорії рандомізованих апроксимацій \cite{briend2023quality} буде проведено ґрунтовні обчислювальні дослідження для оцінки похибок апроксимації на проміжних значеннях глибини; крім цього буде порівняно складність алгоритмів і вдосконалено методи для робастного багатовимірного аналізу.

\begin{thebibliography}{9}
  \bibitem{briend2023quality} S.~Briend, G.~Lugosi, R.I.~Oliveira. \emph{On the quality of randomized approximations of Tukey's depth}. arXiv:2309.05657 (2023).
\end{thebibliography}

\end{document}